\documentclass[12pt]{article}
\usepackage[utf8]{inputenc}
\usepackage[T1]{fontenc}
\usepackage{hyperref}
\title{Urimo Language Resources for K-12 Curriculum Development}
\author{Compiled by Claude Sonnet 5 (claude-sonnet-5) for Andr\'as Kornai}
\date{}
\begin{document}
\maketitle

\section{Introduction}
Urimo (ISO 639-3 urx, also called Yaugiba) is a Torricelli language of Yaugiba village, Turubu Rural LLG, East Sepik Province, Papua New Guinea. About 800 speakers were reported in 2003. It is a different language from Urim (ISO uri).

\section{Primary Educational Resources}
I found no primers, textbooks or school programmes for Urimo.

\section{Linguistic and Reference Materials}
\begin{itemize}
\item Overview of the language: \url{https://en.wikipedia.org/wiki/Urimo_language}
\item Endangered Languages Project entry (lists the alternate name Yaugiba): \url{https://www.endangeredlanguages.com/lang/10575?hl=ru}
\end{itemize}

\section{Religious and Community Resources}
\begin{itemize}
\item Joshua Project people-group profile. It reports about 1,500 people, poor documentation and limited resources: \url{https://joshuaproject.net/index.php/people_groups/15733}
\end{itemize}

\section{Digital and Audio Resources}
None found.

\section{Curriculum Development Context}
Documentation is very thin, and the figures date from 2003. A curriculum would need to begin with community-led documentation. Neighbouring Torricelli languages may offer models for it.
\end{document}